\chapter{Capstone 案例模板与评分 Rubric}

\section{案例导读：Rubric 是提前公开的证据契约}

Rubric 如果只在期末出现，就会被学生理解成拿分技巧；如果从项目第一周就出现，它就是一份证据契约，告诉学生什么材料会被检查、什么判断不能靠印象、什么缺口会让项目暂时不能评分。

本章的重点不是设计一张更复杂的打分表，而是把 Part I 到 Part V 的能力翻译成可观察的交付标准。一个好的 rubric 应该同时帮助学生规划项目、帮助助教校准反馈、帮助教师保护评分的一致性。

\section{案例目标}

第 39 章把 Capstone 项目组织成 project board，第 40 章讨论选题和范围，第 41 章把报告、展示和最终签字收束成 delivery review。本章再往前推进一步：\textbf{把这些材料整理成一套可以发给学生、也可以被教师稳定评分的 case template 和 rubric。}

本章的目标有四点：

\begin{enumerate}
    \item 理解 rubric 不是最后一刻的打分表，而是项目从第一周就应该遵守的证据契约；
    \item 学会把科学问题、数据可复现性、建模评估、结果解释、代码记录、报告展示和 AI 使用披露放到同一个模板里；
    \item 学会区分“可以评分”“小修后评分”“需要较大修订”和“评分前阻塞”四类出口；
    \item 学会生成一份面向学生的 revision checklist，让评分反馈变成下一步行动，而不是只留下一个分数。
\end{enumerate}

\section{Rubric 的作用不是制造分数，而是固定证据标准}

很多课程项目到了期末才第一次看到评分表。这样做有一个问题：学生会把 rubric 当成“最后如何拿分”的清单，而不是“从一开始如何组织证据”的标准。

对 AI 数据分析 Capstone 来说，rubric 至少承担三件事：

\begin{itemize}
    \item 帮学生知道什么叫一个可检查的科学问题；
    \item 帮助教和教师用相同标准判断项目是否可以评分；
    \item 在模型分数很高但证据链不完整时，给出明确的修订路径。
\end{itemize}

因此，本章的核心立场是：\textbf{好的 rubric 不只给总分，还要暴露阻塞项、最低线和修订动作。}

\section{一个极小的 rubric case template 数据集}

配套 notebook 使用 \nolinkurl{capstone_rubric_case_template_demo.csv}。这是一组完全本地、教学用途的\textbf{Capstone 评分样例卡片}。每一行代表一个项目包的评分前状态，包含：

\begin{itemize}
    \item \texttt{case\_id}；
    \item \texttt{project\_id}；
    \item \texttt{theme}；
    \item \texttt{science\_question\_points}；
    \item \texttt{data\_repro\_points}；
    \item \texttt{modeling\_eval\_points}；
    \item \texttt{interpretation\_limits\_points}；
    \item \texttt{code\_git\_points}；
    \item \texttt{report\_presentation\_points}；
    \item \texttt{ai\_use\_integrity\_status}；
    \item \texttt{signoff\_status}；
    \item \texttt{revision\_burden\_score}；
    \item \texttt{reference\_route}。
\end{itemize}

这里的分数采用 100 分制，但不是平均分配。一个建议权重是：

\begin{itemize}
    \item 科学问题清晰度：15 分；
    \item 数据处理与可复现性：20 分；
    \item 模型建立与评估：25 分；
    \item 结果解释与局限性：20 分；
    \item 代码质量与 Git 记录：10 分；
    \item 报告和展示：10 分。
\end{itemize}

请注意，AI 使用披露和 human signoff 在本章中不是额外加分项，而是评分前 gate。原因很简单：如果 AI-use statement 不清楚，或者最终签字边界没有完成，那么项目不能靠模型分数补救。

\section{先看一个危险 baseline：只按总分排序}

最直觉的评分方式是把六个维度加起来，按总分从高到低排序。这个 baseline 很诱人，因为它简单、可解释、也容易写进电子表格。

\begin{examplebox}[只按总分取前三个项目]
\begin{lstlisting}[style=pythonstyle]
top3 = sorted(
    rows,
    key=total_score,
    reverse=True,
)[:3]
\end{lstlisting}
\end{examplebox}

但这个 baseline 有一个危险：它会把“高分但评分前阻塞”的项目混进来。例如，一个项目的数据、模型和展示都不错，但 AI 使用披露缺失；另一个项目模型分数很高，却没有最终 signoff。它们不应该被直接送去评分。

在本章教学数据上，只按总分取前三个项目时，可直接评分的精度是 0.67，并且混入了一个评分前阻塞项目。表~\ref{tab:ch42_total_score_vs_rubric} 展示了结构化 rubric workflow 如何先检查 gate，再判断最低线和修订负担：

\begin{table}[htbp]
\centering
\small
\setlength{\tabcolsep}{4pt}
\begin{tabular}{@{}p{0.26\textwidth}>{\centering\arraybackslash}p{0.17\textwidth}>{\centering\arraybackslash}p{0.17\textwidth}p{0.32\textwidth}@{}}
\toprule
方法 & 可直接评分精度 & 阻塞项混入率 & 教学解读 \\
\midrule
只按总分取前三个项目 & 0.67 & 0.33 & 总分不能替代 AI 使用披露、签字边界和核心维度最低线 \\
结构化 rubric workflow & 1.00 & 0.00 & 先过 gate，再看总分、最低线和修订负担 \\
\bottomrule
\end{tabular}
\caption{本章 notebook 中两个最小评分流程的对照。重点不是让评分变复杂，而是避免把不可评分项目误当成优秀项目。}
\label{tab:ch42_total_score_vs_rubric}
\end{table}

\section{一个可执行的 rubric workflow}

本章 notebook 的 routing 逻辑分成四层。

\textbf{第一层：评分前 gate。} 如果 \texttt{signoff\_status} 不是 \texttt{signed}，或者 \texttt{ai\_use\_integrity\_status} 不是 \texttt{clear}，项目进入 \texttt{blocked\_before\_grading}。这不是惩罚，而是保护学生、教师和项目本身。

\textbf{第二层：核心维度最低线。} 一个项目不能只靠模型分数拿高分。如果科学问题不清楚、数据无法复现、评估缺少 baseline，或者局限性没有讨论，它就需要修订。

\textbf{第三层：总分和修订负担。} 总分仍然重要，但要和 \texttt{revision\_burden\_score} 一起看。一个 80 分项目如果只需要小修，可以进入小修后评分；一个 82 分项目如果核心评估线没过，仍然应该回到修订队列。

\textbf{第四层：反馈包。} route 之后，workflow 不只输出标签，还应该输出学生下一步要做什么。

最小形式可以写成：

\begin{examplebox}[一个最小 rubric routing 逻辑]
\begin{lstlisting}[style=pythonstyle]
if signoff_status != "signed":
    blocked_before_grading
elif ai_use_integrity_status != "clear":
    blocked_before_grading
elif total >= 85 and ready_floors_pass and revision_burden <= 0.25:
    ready_for_grading
elif total >= 75 and minor_floors_pass and revision_burden <= 0.40:
    strong_with_minor_revision
else:
    needs_revision
\end{lstlisting}
\end{examplebox}

\section{Capstone Evidence Package 应该长什么样}

如果要把这个 rubric 直接发给学生，推荐把项目包组织成下面的 Capstone Evidence Package。它不是八张新的独立 record，而是把前五部分已经固定的 Evidence Record、Data Card、Dataset Contract、Model Experiment Record、Trust Statement 和 AI Usage Log 放到同一个提交包里：

\begin{enumerate}
    \item \textbf{Project summary}：一句科学问题、一句话说明为什么这个问题可以用当前数据回答；
    \item \textbf{Data Card / Dataset Contract}：数据来源、许可证、字段、清洗规则、train / validation / test 划分；
    \item \textbf{Baseline evidence}：至少一个非深度学习或简单模型 baseline，以及它的输出和限制；
    \item \textbf{Model Experiment Record}：模型选择、输入输出、训练设置、失败样本；
    \item \textbf{Evaluation and error artifacts}：指标、验证集、测试集、选择效应、过拟合检查和错误案例；
    \item \textbf{Trust Statement}：关键图表、主要 claim、局限性、不确定度和不支持的解释；
    \item \textbf{Reproducibility evidence}：notebook 复跑说明、环境、随机种子、Git 记录；
    \item \textbf{Integrity and disclosure evidence}：AI Usage Log、AI-use statement、引用核查、human signoff 和公开边界。
\end{enumerate}

这组材料的作用，是让 Capstone 从“做了很多事情”变成“每个 claim 都知道证据在哪里”。

\section{把 Evidence Package 接回评分维度}

为了避免 rubric 看起来像另一张独立表格，建议在学生材料中明确写出 Evidence Package 和评分维度的对应关系：

\begin{itemize}
    \item Project summary 对应科学问题清晰度；
    \item Data Card、Dataset Contract 和 reproducibility evidence 共同对应数据处理与可复现性；
    \item Baseline evidence、Model Experiment Record 和 evaluation artifacts 共同对应模型建立与评估；
    \item Trust Statement 对应结果解释、局限性和不确定度；
    \item Reproducibility evidence 同时支撑代码质量、Git 记录和 notebook 复跑；
    \item Integrity and disclosure evidence 不直接加分，而是 AI 使用披露、引用核查、公开边界和 human signoff 的评分前 gate。
\end{itemize}

这样写的目的，是防止学生把“准备材料”和“争取分数”理解成两件事。真正稳定的 rubric 应该让学生从第一周就知道：每多补一张证据卡，都是在补最终评分所需要的证据，而不是在期末临时迎合格式。

\section{配套 notebook 做了什么}

本章 notebook 采用的是一个 teaching-first 的最小设计：

\begin{itemize}
    \item 先读取 rubric case table，并计算六个评分维度的总分；
    \item 用“只按总分排序”做 naive baseline；
    \item 再实现一个带 gate、最低线和修订负担的 rubric workflow；
    \item 输出 ready、minor revision、needs revision 和 blocked 四个队列；
    \item 为每个非 ready 项目生成 revision checklist；
    \item 最后生成一个可以交给学生或助教使用的 case-template prompt。
\end{itemize}

这份 notebook 的重点不是训练一个自动评分器。它的重点是把评分过程变成可审计、可解释、可反馈的项目管理工具。

\section{AI 助手如何帮助你}

在这个案例里，AI 助手最适合帮助你：

\begin{itemize}
    \item 把一个松散项目整理成 Capstone Evidence Package；
    \item 按 rubric 找出缺失证据、低分维度和评分前阻塞项；
    \item 把教师反馈改写成下一周可以执行的 revision checklist；
    \item 检查 AI-use statement 是否具体、诚实、可复核；
    \item 帮助助教统一反馈格式，减少不同评分者之间的漂移。
\end{itemize}

但 AI 助手不应该替学生决定最终分数，也不应该替教师忽略阻塞项。最终评分必须由人类教师或课程团队完成，并且要保留清楚的判断依据。

\section{练习}

\begin{exercisebox}[基础练习]
把教学数据中一个 \texttt{ready\_for\_grading} 项目的 \texttt{signoff\_status} 改成 \texttt{unsigned}。重新运行 notebook，观察它为什么会直接进入评分前阻塞队列。
\end{exercisebox}

\begin{exercisebox}[进阶练习]
新增一条案例：总分高于 85，但 \texttt{modeling\_eval\_points} 只有 14。预测它会落到哪个 route，并说明为什么总分不能替代核心维度最低线。
\end{exercisebox}

\begin{exercisebox}[开放练习]
把你自己的 Capstone 项目整理成 Capstone Evidence Package。然后用本章 rubric 给自己做一次预评分，列出三条最重要的修订动作。
\end{exercisebox}

\section{本章小结}

Capstone rubric 最重要的功能，不是把复杂项目压成一个分数，而是让每个项目在提交、评分和反馈时都有清楚的证据标准。本章把科学问题、数据可复现性、建模评估、解释局限、代码记录、报告展示、AI 使用披露和 human signoff 放进同一个 workflow。

如果学生从第一周就按照这个模板组织项目，最后的评分就不再是一次突然的裁判，而是一次有记录、有依据、也能继续改进的项目复盘。
